% !TEX program = xelatex
\documentclass[11pt]{article}

\usepackage[a4paper,margin=1in]{geometry}
\usepackage{amsmath,amssymb,amsthm,mathtools}
\usepackage{xcolor}
\usepackage{hyperref}

\hypersetup{
  colorlinks=true,
  linkcolor=blue!60!black,
  citecolor=blue!60!black,
  urlcolor=blue!60!black
}

\newcommand{\C}{\mathbb C}
\newcommand{\Z}{\mathbb Z}
\newcommand{\G}{\Gamma}
\newcommand{\slh}{\widehat{\mathfrak{sl}}_2}
\newcommand{\D}{\mathcal D}
\newcommand{\Ind}{\operatorname{Ind}}
\newcommand{\ad}{\operatorname{ad}}
\newcommand{\im}{\operatorname{Im}}

\theoremstyle{plain}
\newtheorem{lemma}{Lemma}
\newtheorem{proposition}[lemma]{Proposition}
\newtheorem{theorem}[lemma]{Theorem}

\theoremstyle{remark}
\newtheorem{remark}[lemma]{Remark}

\title{Boundary Localization of \(M_{-1,1}\), and the General Boundary Version}
\author{Discussion note}
\date{September 2, 2026}

\begin{document}
\maketitle

\section*{0. Setup}

For \(A_1^{(1)}\), let \(\slh=\widehat{\mathfrak{sl}}_2\), with
\[
  [e_i,f_j]=h_{i+j}+i\delta_{i+j,0}K,\qquad
  [h_i,e_j]=2e_{i+j},\qquad
  [h_i,f_j]=-2f_{i+j},\qquad
  [h_i,h_j]=2i\delta_{i+j,0}K.
\]
For integers \(N_1,N_2\), put
\[
  S_{N_1,N_2}
  =
  \operatorname{span}\{K,\ h(n_0),\ e(n_1),\ f(n_2):
  n_0\ge0,\ n_1>N_1,\ n_2>N_2\}.
\]
In particular, \(h_0\in S_{0,0}\). Let
\[
  M=M_{N_1,N_2}(\varphi)=\Ind_{S_{N_1,N_2}}^{\slh}\C_\varphi,
  \qquad
  u=1\otimes1,
\]
and set \(c=\varphi(K)\), \(\lambda_i=\varphi(h_i)\). Then
\[
  e_iu=0\ (i>N_1),\qquad f_ju=0\ (j>N_2),
\]
\[
  Ku=cu,\qquad h_iu=\lambda_iu\ (i\ge0).
\]
Since \(\varphi\) is a character, \([e_i,f_j]=h_{i+j}+i\delta_{i+j,0}K\) with \(i>N_1\),
\(j>N_2\) gives
\[
  \lambda_i=0\qquad(i\ge N_1+N_2+2).
\]

Let \(\G_m\) be the set of finite nondecreasing sequences
\[
  \gamma=(m_1,\ldots,m_k),\qquad m_i\le m,\qquad m_1\le\cdots\le m_k,
\]
and write \(x(\gamma)=x(m_1)\cdots x(m_k)\), \(x(\emptyset)=1\), for \(x=e,h,f\).

\begin{proposition}[PBW basis]\label{prop:pbw}
A PBW basis of \(M_{N_1,N_2}(\varphi)\) is
\[
  \mathcal B_{N_1,N_2}
  =
  \{\,
    e(\alpha)h(\beta)f(\gamma)u:
    \alpha\in\G_{N_1},\ \beta\in\G_{-1},\ \gamma\in\G_{N_2}
  \,\}.
\]
\end{proposition}

\begin{proof}
Take the vector space complement
\[
  \mathfrak a
  =
  \operatorname{span}\{e_i:i\le N_1\}
  \oplus
  \operatorname{span}\{h_i:i\le -1\}
  \oplus
  \operatorname{span}\{f_i:i\le N_2\}
\]
of \(S_{N_1,N_2}\) in \(\slh\). By the PBW theorem,
\[
  U(\slh)\cong U(\mathfrak a)\otimes U(S_{N_1,N_2})
\]
as vector spaces. Tensoring with the one-dimensional \(S_{N_1,N_2}\)-module
\(\C_\varphi\) turns the \(S_{N_1,N_2}\)-factors into scalars, leaving exactly the
ordered PBW monomials of \(\mathfrak a\).
\end{proof}

\begin{lemma}[Localization formulas]\label{lem:loc}
For \(F=f_b\) and \(s\in\Z\),
\[
  f_nF^s=F^sf_n,
\]
\[
  h_nF^s=F^sh_n-2sF^{s-1}f_{n+b},
\]
\[
  e_nF^s
  =
  F^se_n
  +sF^{s-1}\bigl(h_{n+b}+n\delta_{n+b,0}K\bigr)
  -2\binom{s}{2}F^{s-2}f_{n+2b}.
\]
\end{lemma}

\begin{proof}
One has
\[
  [F,f_n]=0,\qquad
  [F,h_n]=2f_{n+b},\qquad
  [F,e_n]=-(h_{n+b}+n\delta_{n+b,0}K),
\]
\[
  (\ad F)^2(e_n)=-2f_{n+2b},\qquad
  (\ad F)^3(e_n)=0.
\]
In the localization,
\[
  xF^s=\sum_{j\ge0}(-1)^j\binom{s}{j}F^{s-j}(\ad F)^j(x),
\]
a finite sum.
\end{proof}

\section*{1. The case \(N_1=-1\), \(N_2=1\)}

Fix
\[
  M=M_{-1,1}(\varphi),\qquad F=f_1.
\]
Since \(F\) is a PBW variable, \(F\) acts injectively on \(M\). Let
\[
  \D_FM=\{F^m:m\ge0\}^{-1}M.
\]
Because the \(f_i\) commute pairwise, \(\D_FM\) has basis
\[
  \{F^s e(\alpha)h(\beta)f(\delta)u:
  s\in\Z,\ \alpha,\beta\in\G_{-1},\ \delta\in\G_0\},
\]
and \(M\) corresponds to \(s\ge0\). The localization quotient
\[
  T_FM=\D_FM/M
\]
has PBW basis
\[
  \mathcal C
  =
  \{F^{-m}e(\alpha)h(\beta)f(\delta)u+M:
  m\ge1,\ \alpha,\beta\in\G_{-1},\ \delta\in\G_0\}.
\]
Put
\[
  \bar u_m=F^{-m}u+M\ (m\ge1),\qquad
  \bar u=\bar u_1.
\]
Define \(\psi:S_{0,0}\to\C\) by
\[
  \psi(K)=c,\qquad
  \psi(h_0)=\lambda_0+2,\qquad
  \psi(h_1)=\lambda_1,\qquad
  \psi(h_i)=0\ (i\ge2),
\]
\[
  \psi(e_i)=0\ (i>0),\qquad
  \psi(f_i)=0\ (i>0),
\]
and let
\[
  A=M_{0,0}(\psi)=\Ind_{S_{0,0}}^{\slh}\C_\psi,
  \qquad
  u'=1\otimes1.
\]

\begin{lemma}\label{lem:cyc}
The vector \(\bar u\in T_FM\) satisfies
\[
  x\bar u=\psi(x)\bar u\qquad(x\in S_{0,0}).
\]
\end{lemma}

\begin{proof}
For \(s=-1\), Lemma~\ref{lem:loc} gives
\[
  f_nF^{-1}=F^{-1}f_n,
\]
\[
  h_nF^{-1}=F^{-1}h_n+2F^{-2}f_{n+1},
\]
\[
  e_nF^{-1}
  =
  F^{-1}e_n
  -F^{-2}\bigl(h_{n+1}+n\delta_{n+1,0}K\bigr)
  -2F^{-3}f_{n+2}.
\]
Hence, for \(n>0\),
\[
  f_n\bar u=0,\qquad e_n\bar u=0
\]
(for \(n=1\), \(f_1\bar u=F^{-1}f_1u+M=F^{-1}Fu+M=u+M=0\)). For the
Cartan part,
\[
  h_0\bar u=(F^{-1}h_0+2F^{-2}f_1)u+M
  =(\lambda_0+2)F^{-1}u+M
  =(\lambda_0+2)\bar u
  =\psi(h_0)\bar u,
\]
\[
  h_1\bar u=(F^{-1}h_1+2F^{-2}f_2)u+M=\lambda_1\bar u=\psi(h_1)\bar u,
\]
and, for \(n\ge2\),
\[
  h_n\bar u=(F^{-1}h_n+2F^{-2}f_{n+1})u+M=0=\psi(h_n)\bar u.
\]
Also \(K\bar u=c\bar u=\psi(K)\bar u\).
\end{proof}

\begin{proposition}\label{prop:phi}
There is a unique \(\slh\)-module homomorphism
\[
  \Phi:A\longrightarrow T_FM,\qquad \Phi(u')=\bar u.
\]
If \(\lambda_1\ne0\), then \(\Phi\) is surjective.
\end{proposition}

\begin{proof}
Let
\[
  \psi_U:U(S_{0,0})\longrightarrow\C
\]
be the algebra character extending the Lie algebra character \(\psi\). By
Lemma~\ref{lem:cyc}, for every \(a\in U(S_{0,0})\) we have
\[
  a\bar u=\psi_U(a)\bar u.
\]
This relation is not only a relation for the generators in \(S_{0,0}\): if
\(a=s_1\cdots s_k\), with \(s_i\in S_{0,0}\), then
\(a\bar u=\psi(s_1)\cdots\psi(s_k)\bar u=\psi_U(a)\bar u\), and the
general case follows by linearity.

Write \(\C_\psi=\C\mathbf 1_\psi\). On the unbalanced tensor product define
\[
  \widetilde\Phi:U(\slh)\otimes\C_\psi\longrightarrow T_FM,
  \qquad
  \widetilde\Phi(x\otimes b\mathbf 1_\psi)=b\,x\bar u.
\]
For \(x\in U(\slh)\), \(a\in U(S_{0,0})\), and \(b\in\C\), the balancing
relation is sent to
\[
\begin{aligned}
  &\widetilde\Phi\bigl(xa\otimes b\mathbf 1_\psi
      -x\otimes a(b\mathbf 1_\psi)\bigr) \\
  &\qquad=b\,xa\bar u-b\,\psi_U(a)x\bar u
   =b\,x\bigl(a\bar u-\psi_U(a)\bar u\bigr)=0.
\end{aligned}
\]
Hence \(\widetilde\Phi\) kills the balancing subspace and descends to the
well-defined linear map
\[
  \Phi:A=U(\slh)\otimes_{U(S_{0,0})}\C_\psi\longrightarrow T_FM,
  \qquad \Phi(x\otimes_{U(S_{0,0})}b\mathbf 1_\psi)=b\,x\bar u.
\]
For \(y\in\slh\),
\[
  \Phi\bigl(y\cdot(x\otimes_{U(S_{0,0})}b\mathbf 1_\psi)\bigr)
  =b\,yx\bar u
  =y\cdot\bigl(b\,x\bar u\bigr),
\]
so \(\Phi\) is an \(\slh\)-module homomorphism. In particular,
\(\Phi(u')=\bar u\); since \(A=U(\slh)u'\), this condition also gives
uniqueness.

For \(m\ge1\),
\[
  e_0F^{-m}=F^{-m}e_0-mF^{-m-1}h_1-m(m+1)F^{-m-2}f_2.
\]
Thus
\[
  e_0\bar u_m=-m\lambda_1\bar u_{m+1},
\]
and
\[
  e_0^r\bar u=(-1)^rr!\lambda_1^r\bar u_{r+1}\qquad(r\ge0).
\]
If \(\lambda_1\ne0\), then
\[
  \bar u_m=\frac{(-1)^{m-1}}{(m-1)!\lambda_1^{m-1}}e_0^{m-1}\bar u\in\im\Phi.
\]
A PBW basis of \(T_FM\) is
\[
  B_{m,\eta}:=F^{-m}X_\eta u+M,
  \qquad
  X_\eta=e(\alpha)h(\beta)f(\delta),
\]
where \(m\ge1\), \(\alpha,\beta\in\G_{-1}\), \(\delta\in\G_0\). Fix a total PBW
monomial order \(\prec\) on the exterior multi-indices \(\eta\), for instance the
degree-lexicographic order induced by the chosen PBW ordering of the generators.
For
\[
  C_{m,\eta}=\frac{(-1)^{m-1}}{(m-1)!\lambda_1^{m-1}}X_\eta e_0^{m-1}u'\in A,
\]
one has
\[
  \Phi(C_{m,\eta})=X_\eta\bar u_m.
\]
Moving the exterior factors in \(X_\eta\) past \(F^{-m}\) gives a finite triangular
expansion
\[
  \Phi(C_{m,\eta})
  =
  B_{m,\eta}
  +
  \sum_{\eta'\prec\eta,\,m'}a_{m',\eta'}B_{m',\eta'}
  +
  \sum_{i=0}^{N_\eta}r_i(m,\eta)\bar u_{m+i},
\]
where \(a_{m',\eta'},r_i(m,\eta)\in\C\); the last finite sum is precisely the part
with zero exterior PBW index.
We prove \(B_{m,\eta}\in\im\Phi\) by induction on \(\eta\). The case \(\eta=0\) is
\(\bar u_m\in\im\Phi\). Assuming
\[
  B_{m',\eta'}\in\im\Phi\qquad(\eta'\prec\eta,\ m'\ge1),
\]
the triangular formula gives
\[
  B_{m,\eta}
  =
  \Phi(C_{m,\eta})
  -
  \sum_{\eta'\prec\eta,\,m'}a_{m',\eta'}B_{m',\eta'}
  -
  \sum_{i=0}^{N_\eta}r_i(m,\eta)\bar u_{m+i}
  \in\im\Phi.
\]
Hence every PBW basis vector of \(T_FM\) lies in \(\im\Phi\).
\end{proof}

\begin{theorem}\label{thm:iso}
Assume
\[
  \lambda_1=\varphi(h_1)\ne0.
\]
Then
\[
  \Phi:M_{0,0}(\psi)\longrightarrow T_{f_1}M_{-1,1}(\varphi)
\]
is an isomorphism.
\end{theorem}

\begin{proof}
By Proposition~\ref{prop:phi}, \(\Phi\) is surjective. For injectivity, take for
\(A=M_{0,0}(\psi)\) the PBW basis
\[
  \mathcal B_{0,0}
  =
  \{\,
    e_0^pe(\alpha)h(\beta)f(\delta)u':
    p\ge0,\ \alpha,\beta\in\G_{-1},\ \delta\in\G_0
  \,\},
\]
and order both \(\mathcal B_{0,0}\) and \(\mathcal C\) lexicographically by
\((p,\eta)\) and \((m,\eta)\), respectively, with \(\eta\) ordered by the
degree-lexicographic order induced by the chosen PBW ordering, under the
correspondence
\[
  e_0^pe(\alpha)h(\beta)f(\delta)u'
  \longleftrightarrow
  F^{-p-1}e(\alpha)h(\beta)f(\delta)u+M.
\]
By Lemma~\ref{lem:loc}, the leading term of
\[
  \Phi(e_0^pe(\alpha)h(\beta)f(\delta)u')
\]
is
\[
  (-1)^pp!\lambda_1^p\,F^{-p-1}e(\alpha)h(\beta)f(\delta)u+M;
\]
the remaining terms come from the commutators produced when moving
\(e_0,h_i,e_i,f_i\) past powers of \(F\), and they either involve lower PBW
monomials or lie in lower filtration layers under the same monomial. Thus
\(\Phi\) is upper triangular between the two PBW bases, with diagonal
coefficients
\[
  (-1)^pp!\lambda_1^p\ne0.
\]
Therefore \(\Phi\) is injective.
\end{proof}

\begin{remark}
If \(\lambda_1=0\), the isomorphism generally fails. Then
\[
  e_0\bar u=0,
\]
while in \(M_{0,0}(\psi)\) the element \(e_0\) is a free variable of the PBW
complement and cannot be killed by the generating vector.
\end{remark}

\begin{theorem}\label{thm:simple}
Assume
\[
  \lambda_1=\varphi(h_1)\ne0,\qquad c=\varphi(K)\ne-2.
\]
Then
\[
  T_{f_1}M_{-1,1}(\varphi)=\D_{f_1}M_{-1,1}(\varphi)/M_{-1,1}(\varphi)
\]
is a simple \(\slh\)-module.
\end{theorem}

\begin{proof}
The FGXZ irreducibility criterion gives simplicity of \(M_{N_1,N_2}(\chi)\)
when
\[
  \chi(h(N_1+N_2+1))\ne0,\qquad \chi(K)\ne-2.
\]
For \(A=M_{0,0}(\psi)\), this is
\[
  \psi(h_1)=\lambda_1\ne0,\qquad \psi(K)=c\ne-2.
\]
Thus \(A\) is simple, and by Theorem~\ref{thm:iso} it is isomorphic to
\(T_{f_1}M_{-1,1}(\varphi)\). Hence the latter is simple.
\end{proof}

\section*{2. The general boundary version}

\begin{theorem}\label{thm:general}
Let \(N\ge1\),
\[
  M=M_{-1,N}(\varphi),\qquad F=f_N,
\]
and assume
\[
  \lambda_N=\varphi(h_N)\ne0.
\]
Then
\[
  \D_{f_N}M_{-1,N}(\varphi)/M_{-1,N}(\varphi)
  \cong
  M_{0,N-1}(\psi),
\]
where \(\psi:S_{0,N-1}\to\C\) is given by
\[
  \psi(K)=c,\qquad
  \psi(h_0)=\lambda_0+2,\qquad
  \psi(h_i)=\lambda_i\ (i\ge1).
\]
If moreover \(c\ne-2\), then the quotient is a simple \(\slh\)-module.
\end{theorem}

\begin{proof}
The proof for \(N=1\) carries over verbatim with the shifts
\(n+1\rightsquigarrow n+N\), \(n+2\rightsquigarrow n+2N\). Lemma~\ref{lem:loc}
gives, for \(s=-1\),
\[
  f_nF^{-1}=F^{-1}f_n,
\]
\[
  h_nF^{-1}=F^{-1}h_n+2F^{-2}f_{n+N},
\]
\[
  e_nF^{-1}
  =
  F^{-1}e_n
  -F^{-2}\bigl(h_{n+N}+n\delta_{n+N,0}K\bigr)
  -2F^{-3}f_{n+2N}.
\]
Hence \(x\bar u=\psi(x)\bar u\) for \(x\in S_{0,N-1}\): for \(n=N\) one uses
\[
  f_N\bar u=F^{-1}f_Nu+M=F^{-1}Fu+M=u+M=0,
\]
the Cartan and \(e\)-remainders are killed by \(\lambda_i=0\ (i>N)\), and
\(h_0\bar u=(\lambda_0+2)\bar u\) since \(2F^{-2}f_Nu=2F^{-2}Fu=2\bar u\).
Moreover
\[
  e_0F^{-m}=F^{-m}e_0-mF^{-m-1}h_N-m(m+1)F^{-m-2}f_{2N},
\]
so
\[
  e_0\bar u_m=-m\lambda_N\bar u_{m+1}
\]
(because \(f_{2N}u=0\)), and \(e_0^r\bar u=(-1)^rr!\lambda_N^r\bar u_{r+1}\).
The PBW bases of the quotient and of \(M_{0,N-1}(\psi)\) are
\[
  \{F^{-m}e(\alpha)h(\beta)f(\delta)u+M:
  m\ge1,\ \alpha,\beta\in\G_{-1},\ \delta\in\G_{N-1}\}
\]
and
\[
  \{e_0^pe(\alpha)h(\beta)f(\delta)u':
  p\ge0,\ \alpha,\beta\in\G_{-1},\ \delta\in\G_{N-1}\},
\]
and the triangular argument of Theorem~\ref{thm:iso} yields the isomorphism,
the diagonal coefficients being
\[
  (-1)^pp!\lambda_N^p\ne0.
\]
Simplicity follows from the FGXZ criterion applied to \(M_{0,N-1}(\psi)\), which
requires
\[
  \psi(h_N)=\lambda_N\ne0,\qquad \psi(K)=c\ne-2.
\]
\end{proof}

\begin{remark}[General \((a,b)\)]
For \(M_{a,b}(\varphi)\) localized at the boundary \(f_b\), assuming
\[
  \varphi(h_{a+b+1})\ne0,
\]
the same triangular PBW argument is expected to give
\[
  \D_{f_b}M_{a,b}(\varphi)/M_{a,b}(\varphi)
  \cong
  M_{a+1,b-1}(\psi),
\]
where \(\psi:S_{a+1,b-1}\to\C\) is given by
\[
  \psi(K)=c,\qquad
  \psi(h_0)=\lambda_0+2,\qquad
  \psi(h_i)=\lambda_i\ (i\ge1),
\]
\[
  \psi(e_i)=0\ (i>a+1),\qquad \psi(f_j)=0\ (j>b-1).
\]
By the FGXZ criterion, \(M_{a+1,b-1}(\psi)\) is simple when
\[
  \psi(h_{a+b+1})=\varphi(h_{a+b+1})\ne0,\qquad \psi(K)=c\ne-2,
\]
in which case the localization quotient is simple.
\end{remark}

\end{document}
